\documentclass[11pt]{article}

\usepackage[utf8]{inputenc}
\usepackage[T1]{fontenc}
\usepackage{lmodern}
\usepackage{geometry}
\usepackage{amsmath,amssymb,amsfonts,amsthm,mathtools}
\usepackage{bm}
\usepackage{enumitem}
\usepackage{microtype}
\usepackage{hyperref}
\usepackage{titlesec}
\usepackage{booktabs}
\usepackage{array}
\usepackage{setspace}
\usepackage{mathrsfs}
\usepackage{physics}

\geometry{margin=1in}
\hypersetup{colorlinks=true, linkcolor=black, urlcolor=blue, citecolor=black}
\setlist[enumerate]{topsep=0.4em,itemsep=0.35em}
\setlist[itemize]{topsep=0.3em,itemsep=0.25em}
\titleformat{\section}{\large\bfseries}{\thesection.}{0.5em}{}
\titleformat{\subsection}{\normalsize\bfseries}{\thesubsection.}{0.5em}{}
\setstretch{1.06}

\newcommand{\Aether}{\AE{}ther}
\newcommand{\Aflow}{\AE{}ther-flow}
\newcommand{\TheoryName}{\Aflow{} Interpretation of Relativity}
\newcommand{\TheoryProgram}{\Aether{} / \Aflow{} framework}
\newcommand{\CompletedMetric}{g^{\sharp}}
\newcommand{\CompletionCore}{\mathfrak C^{\mathrm{zr,aff}}}
\newcommand{\AffineNucleus}{\mathfrak N^{\mathrm{aff}}}

\newtheorem{definition}{Definition}
\newtheorem{proposition}{Proposition}
\newtheorem{remark}{Remark}
\newtheorem{corollary}{Corollary}
\newtheorem{theorem}{Theorem}

\title{\textbf{The \TheoryName{}}\\[0.4em]
\large Primitive-Reservoir Observer-Reduction\\
\large Minimal Zero-Response Affine Completion Core\\
\large Next Benchmark Criterion}
\author{}
\date{}

\begin{document}

\maketitle

\begin{abstract}
The active primitive-reservoir observer-reduction line has now spent the honest bridge-object move on the exact-shell affine nucleus itself. The direct nucleus-forcing theorem already shows that the minimal exact-shell affine nucleus \(\AffineNucleus\) is a canonical output of the explicit minimal zero-response affine completion core \(\CompletionCore\). The remaining honest burden is therefore no longer to restate the nucleus, and it is no longer to reopen the larger primitive reversible constrained dynamics package that the earlier collapse theorem had already screened out.

This note records the routing consequence of that theorem. The next honest benchmark-facing manuscript must therefore do one of two sharper things: either derive or uniquely force the completion core \(\CompletionCore\) from still more primitive explicit substrate structure in a way that changes benchmark-facing status, or prove a genuinely smaller theorem-level collapse, sharper necessity result, or sharper uniqueness result strictly inside \(\CompletionCore\). Another same-output deeper-origin relay that merely preserves the same completion-core consequence does not count next.

The benchmark boundary remains fixed throughout. The one operative metric is still exactly \(\CompletedMetric\), the ordinary matter route is unchanged, there is no second operative metric, the low-energy matter law is not enlarged, and stronger first-principles promotion language remains paused. The positive result of this note is therefore routing discipline at the current completion-core frontier.
\end{abstract}

\section{Introduction}

The active derivational program of \TheoryName{} is no longer at a stage where the exact-shell affine nucleus should be treated as the default unexplained stopping point. The earlier bridge-object compression theorem had already shown that the larger primitive reversible constrained substrate dynamics normal form carries no extra benchmark-facing freedom beyond a smaller minimal exact-shell affine nucleus \(\AffineNucleus\). The later direct forcing theorem then spent the next honest move available there: \(\AffineNucleus\) itself is now forced by the explicit minimal zero-response affine completion core \(\CompletionCore\).

The present note records the routing consequence of that theorem. It does not reopen the completion-core mathematics, and it does not claim a completed first-principles substrate derivation of gravity. It records only which kind of next manuscript would now count as an honest benchmark-facing gain below the current completion-core frontier.

\paragraph{Framework and claim boundary.}
\TheoryProgram{} denotes the broader \Aether{} / \Aflow{} framework built on the claims that the \Aether{} is the underlying four-dimensional substrate of reality, the \Aflow{} is its intrinsic ordered motion, observed three-dimensional space is the local experiential slice of that deeper substrate, S-time is the experienced order of change arising from matter, light, and the \Aflow{}, and observed expansion is the three-dimensional appearance of deeper four-dimensional ordered motion. Within that framework, \TheoryName{} denotes the exact-closure theory in which the effective gravitational dynamics are adopted to be exactly Einsteinian with universal matter coupling. In the active sequence, \emph{adoption} means use of that established relativistic dynamics without claiming substrate derivation, while \emph{derivation} is reserved for a first-principles recovery from explicit substrate variables.

The present note stays strictly inside that boundary. It records only which kind of next manuscript would now count as an honest benchmark-facing gain below the current completion-core frontier.

\section{Fixed Inputs}

\begin{definition}[Completion-core-frontier fixed inputs]
The criterion recorded here uses the following fixed inputs.
\begin{enumerate}
    \item The benchmark exact-closure package, which fixes the public one-metric GR target of the repository.
    \item The benchmark gatekeeping layer, including the post-bridge status correction and the upstream-compression safeguard that blocks same-output deeper-origin relays from counting as new front-line progress once the benchmark-relevant object is unchanged.
    \item The recorded completion-core line, which already compresses the affine-response-graph and combined-response layers to the smaller explicit minimal zero-response affine completion core \(\CompletionCore\).
    \item The direct exact-shell-affine-nucleus forcing theorem, which proves that the minimal exact-shell affine nucleus \(\AffineNucleus\) is canonically forced by \(\CompletionCore\) and therefore removes \(\AffineNucleus\) as an independently unexplained stopping point on the active line.
\end{enumerate}
\end{definition}

\begin{remark}
The present note does not replace those manuscripts. It records the routing consequence of reading them together under the active safeguard against recursive depth drift.
\end{remark}

\section{Why the Completion-Core Forcing Theorem Is the Frontier}

\begin{definition}[Benchmark-neutral completion-core relay]
A \emph{benchmark-neutral completion-core relay} is a manuscript that preserves the same benchmark one-metric output, the same ordinary matter route, the same exact-shell affine nucleus consequence, and the same claim boundary while merely relocating the unexplained substrate layer deeper, restating the nucleus, or re-expanding an older larger shell-facing package, without supplying a new collapse, necessity, uniqueness, or comparably sharp routing gain inside or below the completion core.
\end{definition}

\begin{proposition}[The direct completion-core forcing theorem is the live frontier]
At the current stage of the primitive-reservoir observer-reduction line, the theorem deriving the minimal exact-shell affine nucleus from the minimal zero-response affine completion core remains the live benchmark-facing frontier.
\end{proposition}

\begin{proof}
That theorem changes benchmark-facing status at current scope because it removes \(\AffineNucleus\) itself as the unexplained stopping point. The current bridge object is no longer merely a smaller shell-facing package; it is now a canonical output of the explicit completion core \(\CompletionCore\). By contrast, a manuscript that preserves the same benchmark package and the same completion-core consequence while merely pushing the unexplained layer deeper, rephrasing the nucleus, or reopening the larger shell package does not alter benchmark-facing status unless it adds a comparably sharp gain. Therefore the live frontier is the direct completion-core forcing theorem rather than the earlier collapse theorem or a later same-output relay.
\end{proof}

\begin{corollary}[Status of the nucleus as a next-step target]
Under the current routing safeguard, the exact-shell affine nucleus \(\AffineNucleus\) is no longer itself the default next benchmark-facing target.
\end{corollary}

\begin{proof}
The direct forcing theorem already proves that \(\AffineNucleus\) is forced by \(\CompletionCore\). Once that is on record, another manuscript that spends the move on the nucleus again does not remove any further benchmark-relevant freedom. The open burden has therefore moved below the nucleus itself.
\end{proof}

\section{The Next Benchmark Criterion}

\begin{definition}[Admissible next benchmark-facing completion-core manuscript]
At the current completion-core frontier, a manuscript counts as an admissible next benchmark-facing move only if it does at least one of the following:
\begin{enumerate}
    \item derives or uniquely forces the minimal zero-response affine completion core \(\CompletionCore\) from still more primitive explicit substrate structure in a way that does more than merely relocate the unexplained layer deeper;
    \item proves a genuinely smaller theorem-level collapse, sharper necessity result, or sharper uniqueness result strictly inside \(\CompletionCore\);
    \item does either of the previous two while preserving the same benchmark one-metric package, the same ordinary matter route, and the same claim boundary.
\end{enumerate}
\end{definition}

\begin{theorem}[Next honest benchmark-facing task at completion-core scope]
The next honest benchmark-facing manuscript at the current completion-core frontier must either derive or uniquely force the minimal zero-response affine completion core from still more primitive substrate structure in a benchmark-relevant way, or else prove a genuinely smaller collapse, sharper necessity result, or sharper uniqueness result inside that completion core. A manuscript that only restates the exact-shell affine nucleus, only re-expands the larger primitive reversible constrained dynamics package, or only repeats the same completion-core consequence at a deeper level without such a new gain does not satisfy the criterion.
\end{theorem}

\begin{proof}
By Proposition 1, the live frontier already lies at the theorem that forces the exact-shell affine nucleus from the completion core. Therefore a second manuscript below it must improve on that status rather than merely accompany it. A deeper-origin manuscript can do so only if it changes benchmark-facing status by more than depth alone. If it simply reproduces the same completion-core consequence through another relay, the routing rule classifies it as benchmark neutral. The remaining honest alternatives are therefore exactly those stated in the definition: either a more primitive derivation or uniqueness forcing of \(\CompletionCore\), or a sharper internal collapse, necessity result, or uniqueness result inside \(\CompletionCore\).
\end{proof}

\section{What the Next Move Should Not Be}

\begin{proposition}[Screened next moves]
At the current completion-core frontier, none of the following is the default next benchmark-facing move:
\begin{enumerate}
    \item another restatement of the minimal exact-shell affine nucleus \(\AffineNucleus\);
    \item another re-expansion back to the larger primitive reversible constrained dynamics package;
    \item another same-output deeper-origin relay that merely preserves the same completion-core consequence;
    \item a reopening of older screened layers as if they again defined the live burden;
    \item a manuscript that introduces a second operative metric, enlarges the ordinary matter law, or uses stronger promotion language.
\end{enumerate}
\end{proposition}

\begin{proof}
Items (1)--(3) fail the preceding criterion because they do not directly address the current completion-core burden. They revisit either an already-forced bridge object, an already-screened larger shell package, or the same downstream consequence rather than removing the remaining completion-core freedom or proving a smaller internal core. Item (4) likewise does not directly address the live burden at current scope. It reopens lower or screened layers as if they were still frontier-defining rather than settling the present completion-core task. Item (5) violates the fixed claim boundary of the benchmark package and therefore cannot count as a disciplined next step on the active line. The benchmark package remains the exact one-metric GR package already fixed by adoption, so any admissible next manuscript must preserve that boundary.
\end{proof}

\begin{remark}
This screening statement is not a denial that the older shell-facing manuscripts matter historically. It is a routing statement: they are not the honest way to spend the next benchmark-facing move from the current completion-core frontier.
\end{remark}

\section{Conclusion}

The positive result of this note is routing discipline at the current completion-core frontier. The direct completion-core forcing theorem remains the live benchmark-facing gain because it already removes the exact-shell affine nucleus as the unexplained stopping point on the active line. The nucleus therefore no longer sets the honest next burden, and neither does the larger primitive reversible constrained dynamics package that had already been screened out as a larger same-output presentation.

The scope remains narrow and honest. The benchmark package is unchanged: the one operative metric is still exactly \(\CompletedMetric\), the ordinary matter route is unchanged, there is no second operative metric, and the low-energy matter law is not enlarged. Nothing here upgrades adoption to derivation or licenses stronger promotion language.

The next open burden is therefore clear. The next benchmark-facing manuscript should derive or uniquely force the minimal zero-response affine completion core from still more primitive explicit substrate structure in a way that changes benchmark-facing status, unless the sharper honest gain available instead is a genuinely smaller collapse, sharper necessity result, or sharper uniqueness result strictly inside that completion core. That is the clean next manuscript target at current scope.

\end{document}
